% !TEX root = ../artifacts/main.tex
%
% docs/development.tex
% Development chapter (math-centric revision).

\chapter{Development}
\label{chap:development}

This chapter formulates \textsc{EmotionEngine} as a sequence of
estimation and control problems on residual-stream representations. The
focus is on statistical objects, objective functions, and evaluation
criteria.

Let $\mathcal{C}=\{1,\dots,8\}$ denote the Plutchik categories, let
$h_L(x)\in\mathbb{R}^{d}$ be the last-token residual state at layer $L$,
and let $(x_i^+,x_i^-)$ denote a contrastive pair. The objective is to
estimate directions in $\mathbb{R}^{d}$ that control affective semantics
while minimizing non-affective leakage.

\section{Dataset Construction}

The data objective is to construct samples where emotional signal varies
while lexical-semantic content is approximately matched. Formally, we
seek paired observations drawn from
\begin{equation}
  (x^+,x^-) \sim p(x^+,x^-\mid c),\qquad c\in\mathcal{C},
\end{equation}
with high mutual information between $c$ and
$h_L(x^+)-h_L(x^-)$, and low dependence between pair polarity and
nuisance factors (topic, register, length).

Because standard corpora do not provide this pair structure directly, we
construct it in three stages:
\begin{enumerate}
  \item aggregate corpora to maximize support over
    $(\mathcal{C},\mathrm{PAD})$ jointly;
  \item filter noisy labels to reduce bias in contrastive estimators;
  \item synthesize or mine missing pairs under controlled negative
    sampling.
\end{enumerate}

\subsection{Data Sources, Filtering, and Pair Design}
\label{sec:dev:data:sources}
We combine five English resources (Table~\ref{tab:dev:sources}):
GoEmotions~\cite{demszky2020goemotionsdatasetfinegrainedemotions},
SemEval-2018 E-c~\cite{mohammad2018semeval},
DailyDialog~\cite{li2017dailydialogmanuallylabelledmultiturn},
ISEAR~\cite{scherer1990isear}, and
EmoBank~\cite{buechel2022emobankstudyingimpactannotation}. The first four
provide categorical supervision; EmoBank provides continuous
Valence/Arousal/Dominance targets.

\paragraph{Corpus selection rationale.}
Selection follows three criteria: (i) label coverage over all
$\mathcal{C}$ plus a continuous affect space; (ii) domain diversity to
reduce genre-specific confounding; and (iii) sufficient scale to support
balanced per-class pair construction. EmoBank is essential because it
anchors a supervised linear map into PAD space, complementing discrete
category control.

\begin{table}[t]
\centering
\small
\begin{tabular}{lrl}
\toprule
Source & \# rows & Role \\
\midrule
GoEmotions     &  53\,206 & contrastive (27 fine $\to$ Plutchik) \\
SemEval-2018   &  10\,689 & contrastive (E-c, 11 labels) \\
DailyDialog    & 102\,979 & contrastive (7 codes) \\
ISEAR          &   7\,102 & contrastive (7 emotions) \\
EmoBank        &  10\,061 & V/A/D supervision \\
\midrule
\textbf{Total} & \textbf{184\,037} & --- \\
\bottomrule
\end{tabular}
\caption{Source corpora after unification. Counts are post-loader,
pre-quality-filter.}
\label{tab:dev:sources}
\end{table}

\subsection{Quality Filtering and Contrastive Pair Construction}

Each sample is mapped into the unified label space
$\mathcal{C}\cup{\mathrm{neutral}}$.
To reduce effective label noise, we apply an auxiliary classifier
agreement filter.
Specifically, we use \texttt{j-hartmann/emotion-english-distilroberta-base}
as the external emotion classifier.
Let $y_i$ denote the mapped label and $\hat y_i$ the prediction of an
external emotion classifier.
A sample is retained only if
\begin{align}
\hat y_i &= y_i, \
p_\theta(\hat y_i \mid x_i) &\ge \tau,
\end{align}
with confidence threshold $\tau = 0.5$.
This filtering stage removes samples whose affective label is ambiguous
or inconsistent with the auxiliary classifier, improving the stability
of downstream mean-difference estimators.

After filtering, we construct a total of
$8 \times 500 = 4{,}000$ contrastive pairs.
For each target category $c$, positive examples are drawn from
class-$c$ samples, while negative examples are constructed from a mixture
of three sources:

\begin{enumerate}
    \item empirically mined negatives from opposing emotion categories;
    \item LLM-generated rewrites that alter emotional tone while
    approximately preserving semantic content;
    \item template-based pairs used to maintain coverage when mined
    examples are sparse.
\end{enumerate}

The final mixture proportions are
$0.6/0.3/0.1$ for mined, rewritten, and template-generated pairs,
respectively.
\paragraph{Dataset statistics and limitations.}
\label{sec:dev:data:limits}
Although pair counts are balanced across categories, linguistic
diversity remains imbalanced. Consequently, empirical risk estimates are
conditioned on the constructed pair distribution rather than on a
natural-language prior.

\section{Representation Extraction and Split Protocol}
\label{sec:dev:activations}

\subsection{Layer Choice, Geometry, and Reproducibility}
\label{sec:dev:activations:hooks}
For each pair, we extract last-token residual states
$h_L(x_i^+),h_L(x_i^-)$ at layers
$L\in\{13,16,19,22\}$ (relative depth $0.41$--$0.69$ in a 32-block
stack). Layer selection is treated as model selection over validation
shift-accuracy: layer $16$ is used for category steering and layer $19$
for PAD regression.

\paragraph{Activation geometry.}
\label{sec:dev:activations:storage}
The activation dataset at layer $L$ is represented as
\begin{equation}
  H_L^+\in\mathbb{R}^{n\times d},\qquad H_L^-\in\mathbb{R}^{n\times d},
\end{equation}
with $n=4{,}000$ and $d=4{,}096$. This block-matrix form supports
vectorized estimation for both category means and basis decompositions.

\paragraph{Reproducibility controls.}
\label{sec:dev:activations:repro}
Train/validation/test partitions are stratified by category and seeded.
Hence all reported estimates are deterministic functions of
$(\text{seed},\text{split ratio},\text{layer set})$.

\section{Emotion Representation Space: CAA, Basis, and PAD}
\label{sec:dev:caa}

\subsection{Category Prototypes (CAA)}
\label{sec:dev:caa:formula}
Following contrastive activation addition, for category $c$ and layer
$L$ we estimate
\begin{equation}
\label{eq:caa}
  v_{c, L} \;=\; \frac{1}{|P_c|} \sum_{i \in P_c} h_L^{+}(x_i)
            \; -\; \frac{1}{|N_c|} \sum_{i \in N_c} h_L^{-}(x_i),
\end{equation}
where $P_c,N_c$ are positive and negative index sets for class $c$.
Under $h_L^{\pm}=\mu_{c,L}^{\pm}+\varepsilon^{\pm}$, this estimator is
unbiased for $\mu_{c,L}^+-\mu_{c,L}^-$, with variance scaling as
$O(|P_c|^{-1}+|N_c|^{-1})$.

\paragraph{Per-category prototypes.}
\label{sec:dev:caa:tensor}
Stacking all categories and layers yields
$V\in\mathbb{R}^{8\times 4\times 4096}$. Since
$\|v_{c,L}\|_2$ increases with depth, steering scales are normalized by
median direction norm:
\begin{equation}
  \tilde{v}_{c,L}=\frac{v_{c,L}}{\mathrm{median}_{c'}\|v_{c',L}\|_2}.
\end{equation}
This normalization makes coefficients comparable across $(c,L)$.

\paragraph{Limitations of label-bound vectors.}
\label{sec:dev:caa:limits}
By construction, $\{v_{c,L}\}_{c=1}^{8}$ spans at most an
8-dimensional affective subspace of $\mathbb{R}^{d}$. It cannot encode
emotion-relevant variations orthogonal to category means. Empirically,
raw CAA steering at $\alpha=+2$ yields shift accuracy $0.133$ and
monotonicity correlation $\rho=0.157$, motivating richer
label-independent decompositions.

\subsection{Distributional Bases}
\label{sec:dev:basis}

\paragraph{From category means to per-pair $\Delta$.}
\label{sec:dev:basis:delta}
Instead of collapsing to class means, we preserve pair-level structure:
\begin{equation}
  \Delta \in \mathbb{R}^{|\mathcal{T}|\times d},
  \qquad
  \Delta_i = h_L^{+}(x_i)-h_L^{-}(x_i),
\end{equation}
for train index set $\mathcal{T}$. This retains within-class variance and
cross-class covariance.

\paragraph{Decomposers: NMF / PCA / ICA / Sparse Dict.}
\label{sec:dev:basis:decomposers}
For NMF, we sign-split before factorization:
\begin{equation}
  \tilde{\Delta}=\bigl[\mathrm{relu}(\Delta)\mid\mathrm{relu}(-\Delta)\bigr]
  \in\mathbb{R}^{|\mathcal{T}|\times 2d},
  \qquad
  \tilde{\Delta}\approx HW,
\end{equation}
then recover signed components
$W_{\text{signed}}=W_{\text{pos}}-W_{\text{neg}}\in\mathbb{R}^{k\times d}$.

For comparison, we also fit:
\begin{align}
  &\text{PCA:} &&\min_{W,H}\|\Delta-HW\|_F^2,\quad W^\top W=I,\\
  &\text{ICA:} &&\text{maximize non-Gaussianity of }H,\\
  &\text{Dict:} &&\min_{W,H}\|\Delta-HW\|_F^2+\lambda\|H\|_1.
\end{align}
Each method yields directions $W\in\mathbb{R}^{k\times d}$ and
loadings $H\in\mathbb{R}^{n\times k}$.

\paragraph{$k$-sweep and stability across seeds.}
\label{sec:dev:basis:sweep}
We sweep $k\in\{8,16,32,64\}$ across decomposers and seeds. Stability is
quantified by seed-to-seed component matching under cosine similarity:
\begin{equation}
  s(W^{(a)},W^{(b)})
  =\frac{1}{k}\max_{\pi\in S_k}\sum_{j=1}^{k}
  \cos\bigl(w^{(a)}_j,w^{(b)}_{\pi(j)}\bigr).
\end{equation}
ICA at $k=16$ gives the highest self-consistency in this setup.

\paragraph{Label-independence metrics.}
\label{sec:dev:basis:metrics}
A label-independent basis should avoid collapsing into category
prototypes while retaining controllability. We evaluate:
\begin{enumerate}
  \item mutual information $I(H_j;Y)$ per component,
  \item one-vs-rest linear separability on $H_j$,
  \item dominance ratio
    $\max_c\mathbb{E}[H_j\mid Y=c]/\sum_{c'}|\mathbb{E}[H_j\mid Y=c']|$,
  \item silhouette score of $H$ under class labels.
\end{enumerate}
Low MI and dominance, combined with high steering efficacy, indicate
distributed affective factors beyond category axes.

\subsection{Continuous Affect Mapping}
\label{sec:dev:vad}

\paragraph{PAD regression head.}
\label{sec:dev:vad:head}
At layer $L=19$, PAD mapping is fit via ridge regression:
\begin{equation}
  \min_{W,b}\;\|XW^\top+\mathbf{1}b^\top-T\|_F^2+\lambda\|W\|_F^2,
\end{equation}
where $X\in\mathbb{R}^{n\times d}$ are residual features and
$T\in\mathbb{R}^{n\times 3}$ are PAD targets.
With 6\,402 training and 1\,598 test samples, the resulting
$R^2$ scores are $R^2_V=0.561$, $R^2_A=0.245$, and $R^2_D=0.158$.

\paragraph{Mapping bases $\leftrightarrow$ PAD cube.}
\label{sec:dev:vad:bridge}
Linearity implies that any steering direction
$u\in\mathbb{R}^{d}$ maps to PAD as $Wu\in\mathbb{R}^{3}$. Therefore,
category vectors and basis components admit a shared continuous
interpretation geometry.

\section{Steering and Interaction Model}
\label{sec:dev:steering}

\subsection{Steering Dynamics}
\label{sec:dev:steering:hook}
At selected layers, steering is additive:
\begin{equation}
  h_{L,t} \leftarrow h_{L,t} + \alpha\,m_t\,v,
\end{equation}
where $m_t\in\{0,1\}$ controls token scope (generation-only or all-token)
and $v$ is either a category vector or a basis mixture.

\paragraph{Magnitude scheduling and token scope.}
\label{sec:dev:steering:schedule}
We express magnitude in normalized units,
\begin{equation}
  v=\hat{v}\,\mathrm{median}_{c'}\|v_{c',L}\|_2,
  \qquad \hat{v}=\frac{v}{\|v\|_2},
\end{equation}
and enforce a fluency guardrail by clipping $\alpha$ such that
\begin{equation}
  \frac{\mathrm{PPL}_{\text{steered}}}{\mathrm{PPL}_{\text{base}}}\leq 2.0.
\end{equation}

\paragraph{Steering by basis components.}
\label{sec:dev:steering:basis}
Basis-space control uses linear composition,
\begin{equation}
  v_{\text{mix}}=\sum_{j=1}^{k}\beta_j w_j,
\end{equation}
which extends control trajectories beyond the
8-dimensional category-only subspace.

\subsection{Interactive Interface}
\label{sec:dev:serving}

\paragraph{Interactive control surface.}
\label{sec:dev:serving:api}
The interface implements three mathematical operations: projection,
intervention, and controlled generation. Given input text $x$, analysis
computes residual coordinates $\gamma(x)$; control specifies target
direction $u$ and strength $\alpha$; generation proceeds under the
intervened hidden dynamics.

\paragraph{Dual parameterization (PAD cube and turntable).}
\label{sec:dev:serving:ui}
Two dual parameterizations are exposed: basis coefficients
$\beta\in\mathbb{R}^{k}$ and category weights
$\eta\in\mathbb{R}^{8}$. Their induced steering vector can be written as
\begin{equation}
  v=\sum_{j=1}^{k}\beta_jw_j=\sum_{c=1}^{8}\eta_cv_c,
\end{equation}
allowing direct comparison between distributed and category-bound
control spaces.

\paragraph{Latency and memory considerations.}
\label{sec:dev:serving:latency}
For hidden size $d$, sequence length $S$, and $L_s$ steered layers, the
additional steering-memory term is $O(L_s d)$, typically negligible
relative to model-state memory. Runtime overhead is dominated by
controlled forward passes, trading throughput for intervention fidelity.

\section{Evaluation Protocol and Limits}
\label{sec:dev:practice}

\subsection{Protocol and Reproducibility}
\label{sec:dev:practice:cfg}
All experiments use fixed random seeds, stratified splits, and identical
evaluation prompts across conditions. Hyperparameters are selected on
validation data, and test metrics are reported once per selected model.

Primary evaluation metrics are:
\begin{align}
  &\text{ShiftAcc}=\Pr\bigl[f(\hat{x}_{\alpha})=c^*\bigr],\\
  &\text{Mono}(\alpha)=\rho\bigl(\alpha,s(\hat{x}_{\alpha})\bigr),\\
  &\text{PPLRatio}(\alpha)=
    \frac{\mathrm{PPL}(\hat{x}_{\alpha})}{\mathrm{PPL}(\hat{x}_{0})},
\end{align}
where $f$ is an external affect classifier,
$s$ is a scalar affect score, and
$\hat{x}_{\alpha}$ is generation under steering strength $\alpha$.

\subsection{Validation Scope and Remaining Risks}
\label{sec:dev:practice:tests}
Validation currently emphasizes internal consistency (shape invariants,
determinism, and numerical sanity) and controllability under fixed
decoding settings. Remaining risks are external validity across
domains/languages and partial entanglement between affect and style,
which require broader intervention studies.